% GENERATED FILE. Edit canonical lessons, then run npm run generate:books.
% canonical-source-hash: fnv1a64:eb5f5099a0e8edba
% canonical-lessons: JA-C362-hiyou, JA-C362-keisan-suru, JA-C362-goukei, JA-C362-bukka, JA-C362-shikikin

\chapter{What Will It Cost?}
\label{ch:ja-what-will-it-cost}
\hlchaptermodality{362}

\begin{chapteropening}
\textbf{By the end of this chapter:} \emph{I can say what the cost depends on.}

Complete the last lesson of chapter 362: I can say what the cost depends on.
\end{chapteropening}

\section[\texorpdfstring{\ja{ひよう} (hiyō)}{ひよう (hiyō)}]{\ja{ひよう} (hiyō) — a cost; an expense}

\label{lesson:JA-C362-hiyou}



\begin{quote}
Before the new one: say the Japanese for a loss, then the Japanese for to go out of business.
\end{quote}

\subsection*{You'll want to know: \ja{ひよう}}
\textbf{\ja{ひよう}} — \emph{hiyō} — \textquotedblleft{}a cost; an expense\textquotedblright{}.

\begin{quote}
\textbf{\ja{ひようは} \ja{ばしょに} \ja{よります。}} — \emph{hiyō wa basho ni yorimasu} — \textquotedblleft{}The cost depends on the place.\textquotedblright{}
\end{quote}

\begin{cousinweb}
\textbf{\ja{ひよう}} (\emph{hiyō}) is Sino-Japanese (kango): two Chinese characters read with their on'yomi, readings borrowed from Chinese. They are \emph{hi}, \textquotedblleft{}expense\textquotedblright{} (Middle Chinese \emph{phj+jH}) and \emph{yō}, \textquotedblleft{}utilize\textquotedblright{} (\emph{yowngH}) (KANJIDIC2; Middle Chinese in Baxter's transcription).
\end{cousinweb}

\begin{grammarlens}[title={What you've built}]
One more word for saying what happens if.
\end{grammarlens}

\subsection*{Your turn}
\begin{itemize}
\raggedright
  \item \emph{Say it:} \emph{hiyō}
  \item \emph{Say it:} \emph{hiyō}, once more
  \item \emph{Say it:} \emph{hiyō wa basho ni yorimasu} — The cost depends on the place.
  \item \emph{Recall:} say the Japanese for a loss, then the Japanese for to go out of business, then say \emph{hiyō} again
\end{itemize}

\begin{tcolorbox}[breakable,colback=teal!4,colframe=teal!35!black,title={Before you move on}]
What does \ja{ひよう} mean? (\textquotedblleft{}A cost; an expense\textquotedblright{}.) Say it once more.
\end{tcolorbox}
\section[\texorpdfstring{\ja{けいさんする} (keisan suru)}{けいさんする (keisan suru)}]{\ja{けいさんする} (keisan suru) — to calculate}

\label{lesson:JA-C362-keisan-suru}



\begin{quote}
Before the new one: say the Japanese for to go out of business, then the Japanese for an expense.
\end{quote}

\subsection*{You'll want to know: \ja{けいさんする}}
\textbf{\ja{けいさんする}} — \emph{keisan suru} — \textquotedblleft{}to calculate\textquotedblright{}.

\begin{quote}
\textbf{\ja{かいものが} \ja{おわったら、けいさんして} \ja{ください。}} — \emph{kaimono ga owattara, keisan shite kudasai} — \textquotedblleft{}When the shopping is done, please add it up.\textquotedblright{}
\end{quote}

Polite \textbf{\ja{けいさんします}}, \emph{keisanshimasu}; a \ja{する} verb. Past (polite): \textbf{\ja{てんいんが} \ja{けいさんしました。}}, \emph{ten'in ga keisan shimashita}.

\begin{cousinweb}
\textbf{\ja{けいさんする}} (\emph{keisan suru}) is built from a Sino-Japanese noun and the native verb \emph{suru}, \textquotedblleft{}to do\textquotedblright{}. The noun, \emph{keisan}, is two Chinese characters read with their on'yomi, readings borrowed from Chinese: \emph{kei}, \textquotedblleft{}plan\textquotedblright{} (Middle Chinese \emph{kejH}) and \emph{san}, \textquotedblleft{}calculate\textquotedblright{} (\emph{swanX}) (KANJIDIC2; JMdict; Middle Chinese in Baxter's transcription).
\end{cousinweb}

\begin{grammarlens}[title={What you've built}]
One more word for saying what happens if.
\end{grammarlens}

\subsection*{Your turn}
\begin{itemize}
\raggedright
  \item \emph{Say it:} \emph{keisan suru}
  \item \emph{Say it:} \emph{keisan suru}, once more
  \item \emph{Say it:} \emph{kaimono ga owattara, keisan shite kudasai} — When the shopping is done, please add it up.
  \item \emph{Recall:} say the Japanese for to go out of business, then the Japanese for an expense, then say \emph{keisan suru} again
\end{itemize}

\begin{tcolorbox}[breakable,colback=teal!4,colframe=teal!35!black,title={Before you move on}]
What does \ja{けいさんする} mean? (\textquotedblleft{}To calculate\textquotedblright{}.) Say it once more.
\end{tcolorbox}
\section[\texorpdfstring{\ja{ごうけい} (gōkei)}{ごうけい (gōkei)}]{\ja{ごうけい} (gōkei) — the total; the sum}

\label{lesson:JA-C362-goukei}



\begin{quote}
Before the new one: say the Japanese for an expense, then the Japanese for to calculate.
\end{quote}

\subsection*{You'll want to know: \ja{ごうけい}}
\textbf{\ja{ごうけい}} — \emph{gōkei} — \textquotedblleft{}the total; the sum\textquotedblright{}.

\begin{quote}
\textbf{\ja{ごうけいが} \ja{ちがったら、おしえて} \ja{ください。}} — \emph{gōkei ga chigattara, oshiete kudasai} — \textquotedblleft{}If the total is wrong, please tell me.\textquotedblright{}
\end{quote}

\begin{cousinweb}
\textbf{\ja{ごうけい}} (\emph{gōkei}) is Sino-Japanese (kango): two Chinese characters read with their on'yomi, readings borrowed from Chinese. They are \emph{gō}, \textquotedblleft{}join\textquotedblright{} (Middle Chinese \emph{hop}) and \emph{kei}, \textquotedblleft{}plan\textquotedblright{} (\emph{kejH}) (KANJIDIC2; Middle Chinese in Baxter's transcription).
\end{cousinweb}

\begin{grammarlens}[title={What you've built}]
One more word for saying what happens if.
\end{grammarlens}

\subsection*{Your turn}
\begin{itemize}
\raggedright
  \item \emph{Say it:} \emph{gōkei}
  \item \emph{Say it:} \emph{gōkei}, once more
  \item \emph{Say it:} \emph{gōkei ga chigattara, oshiete kudasai} — If the total is wrong, please tell me.
  \item \emph{Recall:} say the Japanese for an expense, then the Japanese for to calculate, then say \emph{gōkei} again
\end{itemize}

\begin{tcolorbox}[breakable,colback=teal!4,colframe=teal!35!black,title={Before you move on}]
What does \ja{ごうけい} mean? (\textquotedblleft{}The total; the sum\textquotedblright{}.) Say it once more.
\end{tcolorbox}
\section[\texorpdfstring{\ja{ぶっか} (bukka)}{ぶっか (bukka)}]{\ja{ぶっか} (bukka) — prices in general; the cost of living}

\label{lesson:JA-C362-bukka}



\begin{quote}
Before the new one: say the Japanese for to calculate, then the Japanese for the total.
\end{quote}

\subsection*{You'll want to know: \ja{ぶっか}}
\textbf{\ja{ぶっか}} — \emph{bukka} — \textquotedblleft{}prices in general; the cost of living\textquotedblright{}.

\begin{quote}
\textbf{\ja{ぶっかが} \ja{あがったら、こまるでしょう。}} — \emph{bukka ga agattara, komaru deshō} — \textquotedblleft{}If prices go up, people will have a hard time.\textquotedblright{}
\end{quote}

\begin{cousinweb}
\textbf{\ja{ぶっか}} (\emph{bukka}) is Sino-Japanese (kango): two Chinese characters read with their on'yomi, readings borrowed from Chinese. They are \emph{butsu}, \textquotedblleft{}thing\textquotedblright{} (Middle Chinese \emph{mjut}) and \emph{ka}, \textquotedblleft{}value\textquotedblright{} (\emph{kaeH}) (KANJIDIC2; Middle Chinese in Baxter's transcription). The two join as \emph{buk-ka}.
\end{cousinweb}

\begin{grammarlens}[title={What you've built}]
One more word for saying what happens if.
\end{grammarlens}

\subsection*{Your turn}
\begin{itemize}
\raggedright
  \item \emph{Say it:} \emph{bukka}
  \item \emph{Say it:} \emph{bukka}, once more
  \item \emph{Say it:} \emph{bukka ga agattara, komaru deshō} — If prices go up, people will have a hard time.
  \item \emph{Recall:} say the Japanese for to calculate, then the Japanese for the total, then say \emph{bukka} again
\end{itemize}

\begin{tcolorbox}[breakable,colback=teal!4,colframe=teal!35!black,title={Before you move on}]
What does \ja{ぶっか} mean? (\textquotedblleft{}Prices in general; the cost of living\textquotedblright{}.) Say it once more.
\end{tcolorbox}
\section[\texorpdfstring{\ja{しききん} (shikikin)}{しききん (shikikin)}]{\ja{しききん} (shikikin) — a rental deposit}

\label{lesson:JA-C362-shikikin}



\begin{quote}
Before the new one: say the Japanese for the total, then the Japanese for prices in general.
\end{quote}

\subsection*{You'll want to know: \ja{しききん}}
\textbf{\ja{しききん}} — \emph{shikikin} — \textquotedblleft{}a rental deposit\textquotedblright{}.

\begin{quote}
\textbf{\ja{へやを} \ja{でたら、しききんが} \ja{もどるでしょう。}} — \emph{heya o detara, shikikin ga modoru deshō} — \textquotedblleft{}When you move out, you will probably get your deposit back.\textquotedblright{}
\end{quote}

\begin{cousinweb}
\textbf{\ja{しききん}} (\emph{shikikin}) is built from the native \emph{shiki}, \textquotedblleft{}spreading\textquotedblright{} or laying down, a kun'yomi, and the Sino-Japanese \emph{kin}, \textquotedblleft{}gold\textquotedblright{} or money, an on'yomi borrowed from Chinese (Middle Chinese \emph{kim}) (KANJIDIC2; JMdict; Middle Chinese in Baxter's transcription).
\end{cousinweb}

\begin{grammarlens}[title={What you've built}]
One more word for saying what happens if.
\end{grammarlens}

\subsection*{Your turn}
\begin{itemize}
\raggedright
  \item \emph{Say it:} \emph{shikikin}
  \item \emph{Say it:} \emph{shikikin}, once more
  \item \emph{Say it:} \emph{heya o detara, shikikin ga modoru deshō} — When you move out, you will probably get your deposit back.
  \item \emph{Recall:} say the Japanese for the total, then the Japanese for prices in general, then say \emph{shikikin} again
\end{itemize}

\begin{tcolorbox}[breakable,colback=teal!4,colframe=teal!35!black,title={Before you move on}]
What does \ja{しききん} mean? (\textquotedblleft{}A rental deposit\textquotedblright{}.) Say it once more.
\end{tcolorbox}
